% ECONOMICS_PROBLEM: {"id": "C31-121", "title": "Decomposing effects of overlapping mediator mechanisms", "jel_code": "C31", "source_id": "OP-0518"}
% FIRST_POSTED: 2026-10-09
% ATTEMPT_STATUS: unresolved
% ENTRY_KIND: research_attempt
\documentclass[11pt,a4paper]{article}
\usepackage[T1]{fontenc}
\usepackage[utf8]{inputenc}
\usepackage{lmodern}
\usepackage[margin=26mm]{geometry}
\usepackage{amsmath,amssymb,amsthm,mathtools}
\usepackage[hidelinks]{hyperref}
\usepackage{microtype}
\setlength{\parskip}{4pt}
\newtheorem{theorem}{Theorem}[section]
\newtheorem{proposition}[theorem]{Proposition}
\newtheorem{lemma}[theorem]{Lemma}
\newtheorem{corollary}[theorem]{Corollary}
\theoremstyle{definition}
\newtheorem{definition}[theorem]{Definition}
\theoremstyle{remark}
\newtheorem{remark}[theorem]{Remark}
\newcommand{\R}{\mathbb R}
\newcommand{\E}{\mathbb E}
\newcommand{\Prb}{\mathbb P}
\newcommand{\ind}{\mathbf 1}
\newcommand{\Var}{\operatorname{Var}}
\newcommand{\supp}{\operatorname{supp}}
\newcommand{\TV}{\operatorname{TV}}
\DeclareMathOperator*{\argmin}{arg\,min}
\DeclareMathOperator*{\argmax}{arg\,max}
\title{Local mechanism allocations and their statistical error}
\author{GPT 6 Astra Ultra}
\date{9 October 2026}
\begin{document}
\maketitle
\noindent\textbf{Catalogue problem:} \texttt{C31-121}. \textbf{Status:} Unresolved. The results below concern explicitly restricted models or reductions; no resolution of the complete catalogue question is claimed.

\begin{abstract}
With conditionally independent mediators and a sparse low-order outcome expansion, the catalogue's symmetric mechanism allocation is exactly a sum of small local allocations. We prove a polynomial evaluation formula, a simultaneous error bound separating statistical and numerical errors, and an estimable binary subclass. A positive-noise chain shows that bounded parent degree alone does not imply low-order interactions among mechanism switches.
\end{abstract}
\section{The allocation and a restricted graphical model}
Treatment $W\in\{0,1\}$ acts on $p$ mediator mechanisms. For a switched set $S\subseteq[p]$, mechanism $j\in S$ uses its treatment-zero kernel and the others their treatment-one kernels; the outcome mechanism remains that under treatment one. Let $v(S)$ be the resulting mean. The target for mechanism $j$ is
\[
 a_j=\sum_{S\subseteq[p]\setminus\{j\}}
 \frac{|S|!(p-|S|-1)!}{p!}\{v(S\cup\{j\})-v(S)\}.
 \tag{1}
\]
The source motivates representing diffuse and overlapping mediators \cite{CF}; this particular allocation is an editorial choice and is not an intrinsic uniquely determined pathway decomposition.

Assume all causal and positivity restrictions needed to identify the hybrid interventions. In this note's first subclass the mediators are independent conditional on $(X,W)$, each with at most $K$ states, with probabilities $k_j^w(m\mid x)$. Suppose the identified treatment-one outcome mean has a known sparse interaction representation
\[
 \mu_1(x,m)=\sum_{C\in\mathcal C}h_C(x,m_C),
 \qquad |C|\le r,\quad |\mathcal C|=m_0,
 \quad |h_C|\le B_C.
 \tag{2}
\]
An empty interaction is allowed and makes no contribution to (1). Define for $D\subseteq C$
\[
 v_C(D;x)=\sum_{m_C}h_C(x,m_C)
       \prod_{i\in C}k_i^{\ind\{i\notin D\}}(m_i\mid x).
\]
Then $v(S)=\E_X\sum_Cv_C(S\cap C;X)$.
\section{A local identity and exact evaluation}
\begin{theorem}[Local allocation formula]
For $j\in C$, put
\[
 g_{C,j}(x)=\sum_{D\subseteq C\setminus\{j\}}
 \frac{|D|!(|C|-|D|-1)!}{|C|!}
       \{v_C(D\cup\{j\};x)-v_C(D;x)\}.
\]
Then
\[
 a_j=\E_X g_j(X),\qquad
 g_j(x)=\sum_{C\ni j}g_{C,j}(x), \tag{3}
\]
and $\sum_ja_j=v([p])-v(\varnothing)$. Given all nuisance values at one $x$, all $g_j(x)$ can be computed with $O(m_0r2^rK^r)$ arithmetic operations, apart from storage and evaluation of those supplied nuisance functions.
\end{theorem}
\begin{proof}
For a uniformly random permutation of $[p]$, the probability that exactly $S$ precedes $j$ is $|S|!(p-|S|-1)!/p!$. Thus (1) is the expected marginal change on inserting $j$ in that permutation. Expand $v$ using (2). Terms with $j\notin C$ cancel. In a term with $j\in C$, only the relative ordering of the elements of $C$ matters; restriction of a uniform permutation to $C$ is uniform over its $|C|!$ orders. Its predecessor distribution gives precisely $g_{C,j}$. Taking expectation in $X$ proves (3).

For every full permutation, the sum of its $p$ successive marginal changes telescopes to $v([p])-v(\varnothing)$. Averaging proves efficiency. Finally, each of at most $2^{|C|}$ local values requires at most $K^{|C|}$ summands and $O(r)$ operations per product. Summing its marginal changes for all $j\in C$ has no larger order. Add over the $m_0$ interactions.
\end{proof}

\section{A simultaneous error certificate}
Suppose nuisance estimates are trained independently of $n$ evaluation covariates $X_1,\ldots,X_n$. They remain probability kernels, and $|\hat h_C|\le B_C$. On a training event $\mathcal E$ suppose, uniformly in $x$,
\[
 |\hat h_C-h_C|\le\gamma_C,
 \qquad
 \max_{w=0,1}\sum_m|\hat k_i^w(m\mid x)-k_i^w(m\mid x)|\le\delta_i.
 \tag{4}
\]
Let $D_j=\sum_{C\ni j}B_C$ and
$E_C=\gamma_C+B_C\sum_{i\in C}\delta_i$.
Compute local allocations using these fits and average them over the evaluation covariates. Permit an additional deterministic numerical error of absolute size at most $\nu_j$ in the final coordinate; call the output $\hat a_j$.

\begin{theorem}[Uniform coordinate error]
Conditional on any training sample in $\mathcal E$, with probability at least $1-\alpha$ over the evaluation sample, all coordinates satisfy
\[
 |\hat a_j-a_j|\le
 2\sum_{C\ni j}E_C+
 2D_j\sqrt{\frac{2\log(2p/\alpha)}n}+\nu_j.
 \tag{5}
\]
If $\Prb(\mathcal E)\ge1-\eta$, intervals centered at $\hat a_j$ with these radii cover the entire allocation vector with probability at least $1-\alpha-\eta$.
\end{theorem}
\begin{proof}
A product of probability kernels changes in $\ell_1$ distance by at most the sum of the component $\ell_1$ changes. Indeed telescope the two products one factor at a time, sum absolute values, and use unit mass of every unchanged factor. Therefore every local hybrid value has error at most $E_C$. The weights in the local formula for $g_{C,j}$ sum to one, so its two value differences give error at most $2E_C$. Summing proves a uniform approximation error $2\sum_{C\ni j}E_C$ for $g_j$.

Also $|\hat g_{C,j}(x)|\le2B_C$, hence $|\hat g_j(x)|\le2D_j$. Conditional on training, the evaluation variables are independent. The bounded-variable exponential inequality gives
\[
 \Prb\left(\left|n^{-1}\sum_i\hat g_j(X_i)-\E\hat g_j(X)\right|
 >2D_j\sqrt{2\log(2p/\alpha)/n}\right)\le\alpha/p.
\]
A union bound, the deterministic approximation error, and the numerical error give (5). Adding the training failure probability proves the unconditional coverage assertion.
\end{proof}

The formula displays two distinct growth costs: interaction size controls computation, while the number and magnitudes of interactions containing a coordinate control its error. If $B_C\le B$, every coordinate belongs to at most $\Delta$ terms, and all nuisance bounds are $\gamma,\delta$, the first term is at most $2\Delta(\gamma+Br\delta)$.

\section{A directly estimable binary subclass}
There is a useful case in which the local formula simplifies further. Let $M_j\in\{-1,1\}$. Suppose $W$ is randomized with probability $1/2$, and given $X,W=1$ all mediators are independent fair signs. Given $X,W=0$, they are independent with means $\nu_j(x)\in[-1+\zeta,1-\zeta]$, for fixed $\zeta>0$. Assume $0\le Y\le1$ and
\[
 \mu_1(x,m)=h_\varnothing(x)+
          \sum_{\varnothing\ne C\in\mathcal C}h_C(x)\prod_{j\in C}m_j.
 \tag{6}
\]
The joint function is assumed to remain in $[0,1]$. The known collection has $O(p)$ terms, $|C|\le r$, and maximum coordinate incidence $\Delta$. All $h_C,\nu_j$ are $s$-H\"older in $x\in[0,1]^d$, $0<s\le1$, with known radius; $X$ has a density bounded above and below.

\begin{proposition}[Explicit identified formula and constructive rate]
In this subclass,
\[
 h_C(x)=\E\left[Y\prod_{j\in C}M_j\mid X=x,W=1\right],
 \qquad
 a_j=\E_X\sum_{C\ni j}\frac{h_C(X)}{|C|}\prod_{i\in C}\nu_i(X).
 \tag{7}
\]
For fixed $d,r,\Delta$, independent training and evaluation samples of sizes proportional to $n$, histogram estimates of these conditional means with mesh
$h\asymp\{\log(np)/n\}^{1/(2s+d)}$ give simultaneous error
\[
 \max_j|\hat a_j-a_j|
 =O_{\Prb}\left(\Delta\{\log(np)/n\}^{s/(2s+d)}
       +\Delta\sqrt{\log(2p)/n}\right),
 \tag{8}
\]
provided $\log(np)=o(n)$. Known radii and bounded-outcome concentration yield honest simultaneous intervals at the same order. Evaluation is polynomial in $n,p$ for fixed $r,d$.
\end{proposition}
\begin{proof}
The fair independent signs form an orthogonal product basis: the expectation of the product for distinct sets is zero and for the same set is one. Multiplying (6) by the sign product and conditioning proves the first identity. A local interaction has zero mean whenever any member of $C$ remains under its fair treatment-one kernel; if all are switched, its mean is $h_C\prod_{i\in C}\nu_i$. Each member of $C$ is equally likely to be last in a random permutation of $C$, so it receives exactly $1/|C|$ of this value. This proves (7).

For the rate, estimate each $h_C$ by averaging the bounded response $Y\prod_{i\in C}M_i$ among treated observations in a covariate cube, and each $\nu_j$ by averaging $M_j$ among controls there. Clip the fits to $[-1,1]$. Cube approximation bias is at most $L(\sqrt d h)^s$. There are $O(ph^{-d})$ conditional means. Every group--cube probability is at least $c h^d/2$. A binomial lower-tail bound and union bound show that, with probability tending to one, all counts are at least a constant times $nh^d$. Conditional on counts, the bounded-variable inequality and another union bound give a uniform stochastic error at most $C\sqrt{\log(np)/(nh^d)}$. These same inequalities with explicit failure budgets give finite-sample honest bounds, using the full response interval for empty cubes.

Let the resulting uniform bound be $\varepsilon_n$. Since all fitted means have absolute value at most one, telescoping products in (7) bounds the error of its $C$ term by $(1+|C|)\varepsilon_n/|C|\le2\varepsilon_n$. Summing at most $\Delta$ terms gives $2\Delta\varepsilon_n$. The independent evaluation step adds $O_{\Prb}(\Delta\sqrt{\log(2p)/n})$ by bounded-variable concentration. The stated mesh balances bias and stochastic error, proving (8) and its coverage version.
\end{proof}

\section{Why bounded parent sets are insufficient}
For a contrasting identified positive-kernel model, arrange binary mediators in a chain. At the root, $\Prb(M_1=1)$ is $u_0$ or $u_1$ according to its selected mechanism, where $0<u_0,u_1<1$ and $u_0\ne u_1$ (for example, $u_0=1/4$ and $u_1=3/4$). At every subsequent node use
\[
 \Prb(M_j=1\mid M_{j-1}=m)=a+b_wm,
 \qquad w\in\{0,1\},
\]
with $a>0$, distinct $b_0,b_1>0$, and $a+\max b_w<1$. All conditional probabilities are strictly between zero and one. Let the outcome be $Y=M_p$. Every parent set has size one and the outcome is a one-coordinate function. Nevertheless the $p$-fold discrete difference of $v(S)$ with respect to all mechanism switches equals
\[
 (u_0-u_1)(b_0-b_1)^{p-1}\ne0.
\]
To verify this, write recursively $m_j=a+b_{w_j}m_{j-1}$. Terms arising from the additive $a$ do not depend on the root switch and disappear upon differencing that switch. The remaining term is $u_{w_1}\prod_{j=2}^pb_{w_j}$; differencing all coordinates gives the formula. Thus low parent degree and a sparse observed outcome alone do not bound the interaction order of the hybrid value function.

The local algorithm is therefore restricted to independent mediators, not a solution for the full sparse DAG class. Sharp minimax simultaneous rates, general dependence and graph-width conditions, and efficient estimation of arbitrary outcome interactions remain unresolved. Numerical error has been explicitly budgeted rather than folded into an unspecified statistical rate.
\begin{thebibliography}{9}
\bibitem{CF} C. Cinelli, A. Feller, G. Imbens, E. Kennedy, S. Magliacane, and J. Zubizarreta (2025).
\emph{Challenges in Statistics: A Dozen Challenges in Causality and Causal Inference}. arXiv:2508.17099v1, Section 3.4.3, ``Representing complex mediators.''
\url{https://arxiv.org/html/2508.17099v1#S3.SS4.SSS3}.
\end{thebibliography}

\section{Completion Estimate}
30\%

\end{document}
